    %eval-apply
    Interpreters can seem daunting to design, since an interpreter is a program that can carry out any possible computation depending on its input. Yet many share an elegant structure built from two mutually recursive functions: eval, which evaluates expressions in environments, and apply, which applies functions to arguments. They are defined in terms of each other because applying a function means evaluating the expressions in its body, while evaluating an expression may involve applying one or more functions. This back-and-forth is the eval-apply loop, and it continues until the whole program has been reduced to a single value. Gleam’s approach to types and pattern matching suits this design well, since an interpreter must constantly distinguish between the different forms an expression can take.

    %counting calls
    Sometimes it helps to visualize how an interpreter works by counting the calls to eval and apply as we run our code. There are some rules to it. Essentially, every time we encounter a new block of code whose value we're not sure of, we need to evaluate it. In a basic function call, we need to evaluate the function, then evaluate the operands, and only then can we apply the function to its operands. Here's a more detailed example: 
    \begin{lstlisting}[basicstyle=\ttfamily\footnotesize, gobble=8, literate={→}{{$\rightarrow$}}1]
        //Here's the code we're going to run:
        fn f(s) {
            case s {
                [] -> [3]
                [first, ..rest] -> [first, ..f(rest)]
            }
        }
        f([1, 1 + 1])

        //And here's the corresponding series of calls to Eval and Apply:
        eval f([1, 1 + 1])
          eval f
          eval [1, 1 + 1]
            eval 1
            eval 1 + 1
              eval +
              eval 1
              eval 1
              apply + to 1, 1 → 2
          apply f to [1, 2]
            eval case s { [] -> [3] [first, ..rest] -> ... }
              eval s
            eval [first, ..f(rest)]
              eval first
              eval f(rest)
                eval f
                eval rest
                apply f to [2]
                  eval case s { [] -> [3] [first, ..rest] -> ... }
                    eval s
                  eval [first, ..f(rest)]
                    eval first
                    eval f(rest)
                      eval f
                      eval rest
                      apply f to []
                        eval case s { [] -> [3] [first, ..rest] -> ... }
                          eval s
                        eval [3]
                          eval 3
                      apply f to [] → [3]
                apply f to [2] → [2, 3]
          apply f to [1, 2] → [1, 2, 3] //Here it repeats the earlier "apply f to [1,2]", this is not actually an additional apply call, it just restated this apply call once it found the result.
        result: [1, 2, 3]
    \end{lstlisting}
    Something worth noting (also shown as a comment in this trace), is that when we have this gleam interpreter print out the eval/apply trace, if the apply call ends up taking several more steps before finding the result, it will re-print the earlier apply call immediately once it's found the result.
    We can see this with the call to "apply f to [1, 2]", it prints this and then once it's done applying, it prints "apply f to [1, 2] → [1, 2, 3]". We don't count this as an additional call to apply, it just printed this again to demonstrate that it found the result. 
    Conversely, we can see that this call to apply "apply + to 1, 1 → 2" was only printed once, because it immediately found the result and didn't need to have any additional calls to eval/apply before finding the result.
    Ultimately, if we're looking at this type of trace, we have to make sure we disregard any repeated printed calls to apply in order to come to the correct number. 
